% ---------------------------------------------------------------- ঘোষণা
\chapter*{ঘোষণাপত্র}
\addcontentsline{toc}{chapter}{ঘোষণাপত্র}
আমি, \StudentName, এই মর্মে ঘোষণা করছি যে \emph{\ProjectTitle} শীর্ষক এই প্রকল্প প্রতিবেদনটি আমার নিজের কাজ।
ব্যবহৃত সকল তথ্য, ডেটা ও সফটওয়্যারের উৎস যথাযথভাবে স্বীকার ও উল্লেখ করা হয়েছে। এই প্রতিবেদনের সম্পূর্ণ বা
কোনো অংশ অন্য কোনো ডিগ্রি বা যোগ্যতার জন্য দাখিল করা হয়নি।

\vspace{2cm}
\noindent স্বাক্ষর: \rule{6cm}{0.4pt}\hfill তারিখ: \rule{4cm}{0.4pt}

% ---------------------------------------------------------------- সারসংক্ষেপ
\chapter*{সারসংক্ষেপ}
\addcontentsline{toc}{chapter}{সারসংক্ষেপ}
অনেক প্রতিষ্ঠান এখনো বাণিজ্যিক সিদ্ধান্ত নেয় স্প্রেডশিট-ভিত্তিক প্রতিবেদনের উপর নির্ভর করে। এসব প্রতিবেদন অতীতের বর্ণনা দেয়,
কিন্তু খুব কমই তার কারণ ব্যাখ্যা করে বা সামনে কী ঘটতে পারে তা অনুমান করে। এই প্রকল্পে একটি \emph{বুদ্ধিমান ডেটা অ্যানালিটিক্স
ও মেশিন লার্নিং কাঠামো} নকশা, নির্মাণ ও মূল্যায়ন করা হয়েছে। কাঠামোটি ডেটা প্রস্তুতি, অনুসন্ধানমূলক বিশ্লেষণ, পূর্বাভাসমূলক
মডেলিং এবং বিজনেস ইন্টেলিজেন্স ভিজ্যুয়ালাইজেশনকে একটি পুনরুৎপাদনযোগ্য সিদ্ধান্ত-সহায়ক পাইপলাইনে যুক্ত করে। এটি
বাংলাদেশে ড্যান কেক পণ্যের \NRows বিক্রয় লেনদেনের উপর (জানুয়ারি 2020 -- সেপ্টেম্বর 2026) প্রয়োগ করে দেখানো হয়েছে, যেখানে
\NProducts টি পণ্য, \NRegions টি অঞ্চল, \NChannels টি বিক্রয় চ্যানেল ও \NCustTypes ধরনের গ্রাহক রয়েছে।

কাঠামোটি প্রথমে ডেটা যাচাই ও পরিষ্কার করে, একটি ডাইমেনশনাল (স্টার-স্কিমা) ডেটাবেসে সংরক্ষণ করে এবং অনুসন্ধানমূলক ও
অনুমানমূলক পরিসংখ্যান চালায়। এরপর চার ধরনের মেশিন লার্নিং প্রয়োগ করা হয়। \emph{রিগ্রেশন} দিয়ে প্রতিটি পণ্যের মাসিক চাহিদার
পূর্বাভাস দেওয়া হয়েছে; এখানে Linear Regression, Decision Tree, Random Forest, Gradient Boosting, XGBoost,
$k$-Nearest Neighbors ও Support Vector Regression-কে naive ও seasonal-naive বেঞ্চমার্ক এবং Holt--Winters ও SARIMA
মডেলের সঙ্গে তুলনা করা হয়েছে। \emph{ক্লাসিফিকেশন} দিয়ে সাতটি অ্যালগরিদমে বড় অর্ডার পূর্বাভাস করা হয়েছে। \emph{ক্লাস্টারিং}
দিয়ে পণ্য-পোর্টফোলিও ও বাজার বিভাজন করা হয়েছে। \emph{অ্যানোমালি ডিটেকশনে} robust পরিসংখ্যান ও Isolation Forest একত্রে
ব্যবহৃত হয়েছে। সব supervised মডেল 2020--2024 সালের ডেটায় প্রশিক্ষিত এবং পরবর্তী \TestMonths মাসের অদেখা ডেটায় মূল্যায়িত।

ফলাফলে দেখা যায়, নিট বিক্রয় বছরে \CAGRSales\% হারে বেড়েছে, কিন্তু বিক্রির পরিমাণ (volume) বেড়েছে মাত্র \CAGRUnits\%।
অর্থাৎ প্রবৃদ্ধির বড় অংশ এসেছে মূল্য ও পণ্য-মিশ্রণ থেকে, এবং ডিসকাউন্টের কোনো পরিমাপযোগ্য প্রভাব পাওয়া যায়নি। বিক্রয় ঢাকায়
কেন্দ্রীভূত (\DhakaShare\%) এবং ডিসেম্বরে তীব্রভাবে বাড়ে। মাসিক চাহিদা পূর্বাভাসে \BestReg-এর weighted absolute percentage
error ছিল \BestRegWAPE\%, যা seasonal-naive নিয়মের তুলনায় RMSE \GainVsSNaive\% কম। কোম্পানি-পর্যায়ে Holt--Winters মাসিক
ইউনিটের পূর্বাভাস দিয়েছে মাত্র \HWMAPE\% mean absolute percentage error-এ। বড় অর্ডার মাঝারি মাত্রায় পূর্বাভাসযোগ্য
(ROC-AUC \ClsAUC), এবং বিক্রয় চ্যানেলই প্রধান চালিকাশক্তি। ফলাফলগুলো একটি ইন্টারঅ্যাকটিভ ওয়েব ড্যাশবোর্ডে উপস্থাপন
করা হয়েছে এবং স্বয়ংক্রিয়ভাবে প্রমাণ-ভিত্তিক সুপারিশে রূপান্তরিত হয়েছে। গবেষণাটি দেখায় যে একটি সাধারণ, ওপেন-সোর্স প্রযুক্তি
দিয়েই একটি প্রতিষ্ঠান বর্ণনামূলক প্রতিবেদন থেকে পূর্বাভাসমূলক, ডেটা-নির্ভর সিদ্ধান্ত গ্রহণে উত্তরণ করতে পারে।

\vspace{4mm}
\noindent\textbf{মূলশব্দ:} ডেটা অ্যানালিটিক্স; মেশিন লার্নিং; চাহিদা পূর্বাভাস; বিজনেস ইন্টেলিজেন্স; সিদ্ধান্ত-সহায়ক
ব্যবস্থা; গ্রাহক বিভাজন; অ্যানোমালি ডিটেকশন; Python।

% ---------------------------------------------------------------- কৃতজ্ঞতা
\chapter*{কৃতজ্ঞতা স্বীকার}
\addcontentsline{toc}{chapter}{কৃতজ্ঞতা স্বীকার}
প্রকল্পের পুরো সময় জুড়ে দিকনির্দেশনার জন্য আমার তত্ত্বাবধায়ক \Supervisor-এর প্রতি আন্তরিক কৃতজ্ঞতা জানাই। যে জ্ঞানের উপর
এই কাজ দাঁড়িয়ে আছে, তার জন্য \Programme{} প্রোগ্রামের শিক্ষকমণ্ডলীকে ধন্যবাদ। ডেটাসেট সংগ্রহ ও প্রস্তুতিতে যাঁরা সাহায্য
করেছেন এবং ধৈর্য ও সমর্থনের জন্য আমার পরিবারকেও কৃতজ্ঞতা জানাই।
